% Round 18. ci_global branch. Candidate; independent review required.
\section{Iterated complete intersections and a final deep degeneration}
\label{r18ce:section}

Two distinct constructions are kept separate here.  The first takes
sections of an ambient model whose bad fibers have only finitely
many singular points; it produces a new model with good or ordinary
quadratic reduction at every place and can be iterated.  The second
inserts CA18's deep product degeneration only at the last step.
A deep normal-crossings ambient fiber has a positive-dimensional
double locus, so it is not an input to the first construction.
An ample section cannot in general avoid that locus.

\subsection{Native ambient data and the two geometric lemmas}

Let $F$ be a number field and let
\[
       \mathcal Y\subset\mathbf P^N_{\mathcal O_F}
 \tag{CE18.1}\label{r18ce:ambient}
\]
be a specified integral flat projective relative complete
intersection of relative dimension $m\ge4$, cut out by actual
homogeneous equations of fixed positive degrees
$b_1,\ldots,b_c$, where $c=N-m$.
Assume that its generic fiber $Y/F$ is smooth and geometrically
integral, and that every geometric special fiber is integral
and is smooth except at finitely many points.
Good or ordinary-quadratic special fibers satisfy the latter
condition, but the ambient local singularities do not need to
be resolved for the argument below.  The equations and the
relative complete-intersection property are part of the input.
In particular, this is not a claim that an arbitrary variety
admits such a model.

Fix $d\ge3$ and let $V_{d,F}$ be the full vector space of
degree-$d$ homogeneous forms on the displayed projective space.
All forms below are actual forms in this space, even if their
restriction map to $Y$ has a kernel.
The map onto $H^0(Y,\mathcal O_Y(d))$ is surjective: the
Koszul resolution of the complete-intersection ideal gives
$H^1(\mathbf P^N,\mathcal I_Y(d))=0$ in the present
positive-dimensional range.  Thus the section-space density
below has its usual full-linear-system meaning.

\begin{lemma}[Finite singular loci can be avoided at a special fiber]
\label{r18ce:residue-Bertini}
Let $Y_0\subset\mathbf P^N_k$ be a geometrically integral
complete intersection of dimension $m\ge4$ whose geometric
singular locus is finite.  There is a nonempty geometric open
of degree-$d$ forms whose divisor on $Y_0$ is smooth and
geometrically integral and avoids that singular locus.
Over a finite $k$ this open has a point over some finite
extension of $k$; no fixed bound or degree-one assertion
about that extension is required.
\end{lemma}

\begin{proof}
Over $\bar k$, avoiding the finite singular set excludes
finitely many proper linear hyperplanes in coefficient space.
On the smooth locus of $Y_0$, evaluation of first jets is
surjective.  Indeed the monomials $X_0^d,X_0^{d-1}X_i$
give all first jets on projective space at the first
coordinate point, and the quotient to the tangent space of
$Y_0$ is surjective.  The incidence of a singular divisor
point in that smooth locus therefore has dimension at most
one less than the coefficient parameter space.  Its image
closure is a proper subset.  Intersecting the complement
with the open avoiding the finite singular locus gives
a nonempty open of globally smooth divisors.

Each such nonzero divisor is a positive-dimensional
projective complete intersection and is geometrically
connected.  For a direct cohomological check, the Koszul
resolution by the ambient defining equations, together
with the vanishing of intermediate cohomology of line
bundles on projective space, gives
$H^0(Y_0,\mathcal O_{Y_0})=k$ and
$H^1(Y_0,\mathcal O_{Y_0}(-d))=0$ after algebraic closure.
In the resolution of $\mathcal O_{Y_0}(-d)$ every term
is a negative line bundle.  Its only possible nonzero
projective-space cohomology is in degree $N$, and the
resolution has length $c=N-m$; thus no cohomology in
degrees $0<i<m$ survives.  The divisor exact sequence
then gives its $H^0$ equal to $\bar k$.
Smoothness and connectedness imply geometric integrality.
Finally a point in a nonempty open over $\bar k$ has
finitely many algebraic coordinates.  If $k$ is finite,
these coordinates lie in one finite extension.
\end{proof}

\begin{lemma}[Fixed-member pencils on the actual smooth generic ambient]
\label{r18ce:fixed-pencil}
Let $H_0\in V_{d,F}$ cut a smooth divisor on $Y$.
There is a nonempty Zariski open of $H_1\in V_{d,F}$
such that $H_1$ and the base $H_0=H_1=0$ cut smoothly
on $Y$, and every singular geometric pencil member has
one ordinary quadratic point with transverse discriminant
intersection.
\end{lemma}

\begin{proof}
The full proof of CG18's fixed-member lemma applies with
$5,6,5$ replaced by $d,N,m$, respectively; here are the
precise inputs so that this is not an invocation beyond
that proof's checked range.
First jets at one point are supplied by
$X_0^d,X_0^{d-1}X_i$.  Second jets are supplied also by
$X_0^{d-2}X_iX_j$, and are therefore arbitrary on the
smooth tangent space of $Y$.  First jets at two distinct
points are jointly independent: after making those two
coordinate points, the two lists
$X_0^d,X_0^{d-1}X_i$ and
$X_1^d,X_1^{d-1}X_j$ are distinct and each has zero
first jet at the opposite point.  The inequality $d\ge3$
is used exactly here.

If the coefficient projective space has dimension $M$,
the marked-singularity incidence over $Y$ is a smooth
irreducible projective bundle of dimension $M-1$.
The degenerate-Hessian locus and the closure of the
two-marked-point incidence after forgetting its second
point each have dimension at most $M-2$.  Their
complement gives unique nondegenerate singularities;
the coefficient image is an irreducible discriminant
hypersurface with a dense smooth one-node open.
The Hessian implicit-function argument identifies its
normal differential with variation of the critical value.

Projection of that discriminant from the prescribed
smooth point $[H_0]$ is finite, because no line through
that point is contained in the discriminant.  It is
generically etale in characteristic zero.  Exclude its
branch locus and the image of its non-one-node subset.
This gives a nonempty open of transverse one-node
directions.  The smooth-base condition gives another
nonempty open by Bertini on the fixed smooth $H_0$
divisor; the restricted degree-$d$ system is very ample.
Their intersection in the irreducible direction space,
followed by exclusion of singular $H_1$, proves the lemma.
All constructions are defined over $F$.
\end{proof}

\subsection{A section construction with no deep special fiber}

\begin{theorem}[One fixed finite extension and a dense ordinary-reduction family]
\label{r18ce:ordinary-family}
Under \eqref{r18ce:ambient}, for every $d\ge3$ there is a
finite extension $E/F$ and a Zariski-dense set of smooth
divisors $X\subset Y_E$ cut by degree-$d$ forms with a
flat integral projective complete-intersection model over
$\mathcal O_E$ such that every special fiber is
geometrically integral and is either smooth or has one
ordinary quadratic point in odd residue characteristic.
Every place of residue characteristic two is a good
place for $X$.

Every such $X$ has full ordinary actual strong G in
all codimensions at every finite completion, for all
new cycles after every finite extension of $E$.
Consequently PH16 gives the full complementary
geometrically homological rational Chow-group
Beilinson--Bloch heights over every such extension.
The density is in the section space of this one fixed
$Y_E$ and this one fixed $E$.
\end{theorem}

\begin{proof}
\emph{1. Fix a smooth reference member and a finite bad set.}
Choose an integral actual form $H_0$ cutting a smooth
geometrically integral divisor on $Y$, by characteristic-zero
Bertini and clearing denominators.
Let $S_0$ be a finite set of rational primes containing
$2$, every rational prime below a nonsmooth special fiber
of $\mathcal Y$, and every prime below a fiber on which
$H_0$ fails to cut a smooth divisor of dimension $m-1$.
These sets are finite, by properness and the given
smooth generic fibers.  In particular, outside $S_0$
both $\mathcal Y$ and its reference divisor are smooth.
Include all primes of $F$ above the rational primes
in $S_0$ in the next step, including those at which
no failure occurred.

\emph{2. Choose one extension for all residue conditions.}
For each such prime ideal $\mathfrak p$, apply
Lemma~\ref{r18ce:residue-Bertini} to its actual special
fiber.  Choose a smooth section over some finite
extension of $\kappa(\mathfrak p)$, and let its degree
be $f_{\mathfrak p}$.  Choose $N_0\ge2$ divisible by
all these finitely many degrees.
For each $\mathfrak p$ choose a monic irreducible
degree-$N_0$ polynomial over $\kappa(\mathfrak p)$.
The Chinese remainder theorem on coefficients gives
one monic $g\in\mathcal O_F[T]$ with those reductions.
It is irreducible over $F$: a factorization into
monic factors over $F$ would have integral
coefficients, since $\mathcal O_F$ is integrally
closed, and would contradict any of the irreducible
reductions.  Let $E=F[T]/(g)$.
At each of the selected $\mathfrak p$ there is one
prime $\mathfrak P$ of $E$, with unramified local
degree $N_0$ and residue extension of that degree.
This follows either from the finite etale local
algebra $\mathcal O_{F_{\mathfrak p}}[T]/(g)$ or
from the unit discriminant and irreducible reduction.
All the chosen smooth residue sections therefore
exist over their native $\kappa(\mathfrak P)$.
Choose actual integral lifts $G_{\mathfrak P}$ of
their coefficient tuples.  This fixes $E$ once and
for all before varying the next section.

\emph{3. Choose a fixed-member pencil in the coefficient lattice.}
Let $B=\prod_{p\in S_0}p$.  Require
\[
       H_1\equiv G_{\mathfrak P}\pmod{\mathfrak P}
          \quad\text{at every selected }\mathfrak P.
 \tag{CE18.2}\label{r18ce:coefficient-CRT}
\]
The moduli are distinct maximal ideals.  Their solution
set is a nonempty affine lattice in integral form
coefficients over $E$, with a nonzero ideal as its
translation module.  Such an ideal contains a
nonzero integer, so the lattice contains a product
of infinite arithmetic progressions and is Zariski
dense over $E$.  Intersect it with the nonempty open
of Lemma~\ref{r18ce:fixed-pencil}.  The resulting
allowed $H_1$ remain Zariski dense.

After fixing an allowed $H_1$, spread its pencil over
$\mathcal O_E[1/S]$, for a finite set $S$ of prime
ideals containing all primes above $S_0$.  Outside
$S$, the ambient and base are smooth, the pencil
total space is smooth, the relative critical scheme
is finite etale with distinct critical values, and
its Hessians are invertible.  Its map to the
parameter line is a closed immersion.  Exclude
also any prime where the two sections become
dependent.  These are the finite-presentation
spreading arguments proved in CG18: first remove
vertical critical components, then spread all the
listed open conditions and finite morphisms.
Every geometric member outside $S$ is smooth or
has precisely one ordinary quadratic singularity.

\emph{4. Choose the arithmetic parameters and check every place.}
Let $T$ be the finite set of rational primes under
$S$.  Take integers $a$ in the infinite progression
\[
       a\equiv1\pmod B,
       \qquad a\equiv0\pmod p\quad(p\in T\setminus S_0),
 \tag{CE18.3}\label{r18ce:parameter-CRT}
\]
excluding the finitely many singular generic members,
and let
\[
                         X_a=Y_E\cap V(BH_0+aH_1).
 \tag{CE18.4}\label{r18ce:ordinary-equation}
\]
At every prime above $S_0$, its reduction is the
chosen smooth $H_1$ divisor, avoiding every singular
point of the ambient special fiber.  At a prime
above $T\setminus S_0$, $B$ is a unit and the
reduction is the smooth $H_0$ divisor.  At every
remaining prime, $[B:a]$ is an integral parameter
of the spread pencil, so its special fiber is
smooth or has one ordinary quadratic point.
The characteristic-two primes are all in the
first, smooth case.

All these sections are nonzero on the corresponding
integral ambient fibers.  Their divisor sequences
are consequently flat over the native valuation
rings, and they glue to the indicated global
relative complete-intersection model.  The smooth
integral generic fiber and flatness imply that
the model is integral.  Its special fibers with
isolated singularities are geometrically integral:
the complete-intersection condition is $S_2$,
isolated singularities in dimension $m-1\ge3$
give $R_1$, and the Koszul connectedness argument
in Lemma~\ref{r18ce:residue-Bertini} gives
connectedness.  They are thus connected normal
schemes.  At the selected bad-ambient primes,
the divisor is actually contained in the smooth
locus of the ambient over the valuation ring
near the special fiber, so its asserted good
reduction is literal.

\emph{5. Complete cycle domains, extension stability and density.}
The frozen good-reduction theorem and ON17 apply
at every finite place.  ON17 is a theorem for
flat integral projective models of ordinary
quadratic reduction, independent of the number
of global ambient defining equations.  Its
arbitrary-thickness and finite-extension
conclusions apply anew to all cycles over
every finite local extension.  Thus the
full ordinary-actual G domain is obtained,
with exact specified supports and full tests.
PH16 and GA17 give the stated complete height
domains over every finite number-field
extension of $E$.

For each allowed $H_1$ infinitely many integers
$a$ remain.  A polynomial vanishing on all
$BH_0+aH_1$ therefore vanishes on each full
affine line and hence at $BH_0+H_1$.
The allowed $H_1$ form a dense set, so that
polynomial is zero.  Passing to the linear
section-space image and its projectivization
proves the precise density assertion.  No
uniform spreading set as $H_1$ varies is
needed.
\end{proof}

\subsection{Iteration without inserting a deep ambient fiber}

\begin{corollary}[Specified multidegrees]
\label{r18ce:iteration}
Fix a number field $F_{\mathrm{base}}$, a final dimension
$n\ge3$, and a finite multidegree
$(d_1,\ldots,d_t)$ with each $d_i\ge3$.
After a finite extension $E/F_{\mathrm{base}}$ there
are smooth geometrically integral complete intersections
of that multidegree in $\mathbf P_E^{n+t}$ having
good or ordinary quadratic reduction at all finite
places, with good reduction at residue characteristic
two.  They have full local G and the complete
complementary homological Chow-group heights over
every finite extension of $E$.

More precisely one can fix the first $t-1$ defining
equations and one final field $E$ so that a Zariski-dense
set of last degree-$d_t$ sections on that fixed
ambient complete intersection has these properties.
For $t=1$ this reduces to a weaker-field-version of
the already frozen GP17 sector.
\end{corollary}

\begin{proof}
Start with $\mathbf P^{n+t}_{\mathcal O_{F_{\mathrm{base}}}}$.
Apply Theorem~\ref{r18ce:ordinary-family} in degree
$d_1$, and choose one of its smooth generic members
and its actual global model.  That model is a
relative complete intersection with geometrically
integral good or isolated-node special fibers.
It is therefore an input to the same theorem in
degree $d_2$.  Continue in this manner.
Before each cut the ambient relative dimension
is at least $n+1\ge4$, exactly as required.
Each finite field extension preserves the geometric
integrality and isolated-node property of the
previous ambient model; possible node thickness
changes do not alter this special-fiber condition.

There are finitely many cuts, so their fields
form a finite tower with finite total extension
$E/F_{\mathrm{base}}$.  At the last cut keep
its field, its already selected ambient, and
all its residue data fixed, and vary the
last section as in the density theorem.
This proves the stated density and cycle
coverage.  It does not claim joint density
in the space where all prior ambient equations
are simultaneously varied.
\end{proof}

The same iteration can begin with a specified smooth
relative split quadric over $\mathbf Z$, as in CG18,
followed by degrees at least three.  It therefore
also supplies multidegrees $(2,d_2,\ldots,d_t)$ with
$d_i\ge3$ for $i\ge2$.  No assertion about arbitrary
degree-two steps inside an already chosen ambient
is inferred from the degree-$d\ge3$ two-jet proof.

\subsection{A final deep product degeneration}

Now let the input \eqref{r18ce:ambient} have
$m=n+1\ge5$.  This includes any of the ordinary
complete-intersection ambient models just constructed.
Fix $s\ge3$ and positive integers $e_1,\ldots,e_s$
of sum $D$.  Put $\delta=\prod_{j=1}^c b_j$,
with $\delta=1$ for projective-space ambient.

\begin{theorem}[Deep full-height sections of a fixed complete intersection]
\label{r18ce:deep-family}
There is one finite extension $E/F$, one rational
prime $q$ unramified in $E$, and a Zariski-dense
set of smooth degree-$D$ divisors $X\subset Y_E$
such that:
\begin{enumerate}
\item at every finite completion, every specified
      geometrically homologically trivial rational
      cycle has ordinary actual strong G in every
      codimension, with exact support and full tests;
\item at every place above $q$ the native model is
      the CA18 product model with factor degrees
      $e_1,\ldots,e_s$, and
      \[
       \operatorname{rank}N^n
        =\delta\sum_{|I|=n+1}\prod_{i\in I}e_i
                                -\binom{s-1}{n};
       \tag{CE18.5}\label{r18ce:monodromy}
      \]
\item after the independently reviewed RI18 theorem
      is applied at these places, the full assertion
      holds for all new cycles over every finite
      number-field extension of $E$, and PH16 gives
      the complete complementary Chow-group heights
      over every such extension.
\end{enumerate}
At the places away from $q$ the reductions are good
or ordinary quadratic in odd residue characteristic.
The fixed field and fixed ambient are part of the
density assertion.
\end{theorem}

\begin{proof}
Run steps 1 and 2 of
Theorem~\ref{r18ce:ordinary-family} in degree $D$,
fixing $H_0,S_0,E$ and all smooth residue sections.
Before selecting $H_1$, choose actual degree-$e_i$
factor forms over $E$ whose divisors on the smooth
generic complete intersection $Y_E$ form a geometric
SNC divisor with all intersections smooth of the
expected dimension and with no $(n+2)$-fold point.
This is obtained successively by characteristic-zero
Bertini on the finite collection of previous strata
and by avoiding their zero-dimensional intersections.
The positive-dimensional complete intersections are
geometrically connected, hence integral.  Choose
a degree-$D$ form $Q$ transverse to each of those
positive-dimensional strata and avoiding the
zero-dimensional ones.  All these open conditions
have $E$-rational choices because $E$ is infinite.
Clear denominators and put $P=\prod_i f_i$.

The ambient model is smooth away from finitely
many primes, and all the chosen geometric
conditions spread away from finitely many more.
Choose a rational $q$ outside their underlying
finite set, outside $S_0$, and outside the
discriminant of $E/\mathbf Q$.  Then every
$w\mid q$ has a smooth complete-intersection
ambient model and the exact native CA18 data.

Let $B=\prod_{p\in S_0}p$, $M=q^2B$, and add
to the residue congruences \eqref{r18ce:coefficient-CRT}
the congruence
\[
                         H_1\equiv P+qQ\pmod{q^2\mathcal O_E}.
 \tag{CE18.6}\label{r18ce:deep-CRT}
\]
The ideals are coprime, so this is still a full
affine coefficient lattice.  Choose $H_1$ in
its intersection with the fixed-$H_0$ Lefschetz
open.  Spread only after this choice, add all
places above $q$ to the finite exceptional set,
and choose integers $a$ with
\[
 \begin{gathered}
 a\equiv1\pmod{q^2},\qquad a\equiv1\pmod B,\\
 a\equiv0\pmod p
   \quad\text{for the other rational primes under that set}.
 \end{gathered}
 \tag{CE18.7}\label{r18ce:deep-parameters}
\]
Exclude the finitely many singular generic members
and put $X_a=Y_E\cap V(MH_0+aH_1)$.
At $w\mid q$ this is exactly
\[
                  P+qQ_a=0,\qquad
            Q_a=(MH_0+aH_1-P)/q,\qquad\bar Q_a=\bar Q.
\]
The coefficient congruences make the quotient
integral, and $q$ has native uniformizer exponent
one.  CA18 supplies full local G and
\eqref{r18ce:monodromy}.

At every other finite place the proof in step 4
of Theorem~\ref{r18ce:ordinary-family} is unchanged:
the original finitely many bad ambient fibers
are cut by chosen smooth divisors avoiding their
finite singular sets, the other exceptional
places reduce to the smooth $H_0$, and all
remaining members have good or ordinary-node
reduction.  These give full local G.

CA18 applies afresh after every unramified
extension at a place above $q$.  These are
exactly the inputs of RI18's norm-detection
theorem, which covers all cycles over arbitrary
finite extensions, retaining the full field
degree in its final actual trace and the
original absolute value.  Good reduction and
ON17 already have that extension stability
at every other place.  Hence PH16 and GA17
give the stated complete heights over every
finite extension of $E$.

For density use the same line argument with
$MH_0+aH_1$ as in the preceding theorem.
The additional congruence is fixed, so its
full affine lattice remains dense; the prime,
field and ambient are not varied in that
argument.  This proves the theorem.
\end{proof}

If $s\ge n+2$, the integer in
\eqref{r18ce:monodromy} is positive, since it is at
least $\binom{s}{n+1}-\binom{s-1}{n}
=\binom{s-1}{n+1}>0$.
If $s=n+1$, it equals $\delta\prod_i e_i-1$
and is positive exactly when $\delta\prod_i e_i>1$.
In particular, for an ambient complete intersection
of degree $\delta\ge2$, every last degree
$D\ge n+1$ admits such a positive-rank construction
by taking $s=n+1$ and factor degrees
$(1,\ldots,1,D-n)$.
The frozen NM17 obstruction then excludes potential
good reduction and every strict semistable dominant
witness of crossing depth less than $n$.

\subsection{An intrinsic geometric distinction in all these dimensions}

\begin{proposition}[At least two nonlinear complete-intersection equations]
\label{r18ce:not-hypersurface}
Let $X$ be a smooth geometrically integral complete
intersection of dimension $n\ge3$ in projective
space, with multidegree $(a_1,\ldots,a_t)$.
If at least two of the $a_i$ are at least two,
then $X$ is not geometrically isomorphic to any
smooth hypersurface of dimension $n$ in
$\mathbf P^{n+1}$.
\end{proposition}

\begin{proof}
Work over an algebraic closure.  The primary
Grothendieck--Lefschetz theorem gives
\[
                    \operatorname{Pic}(X)=\mathbf Z[\mathcal O_X(1)]
\]
for a projective complete intersection of
dimension at least three, and the same assertion
holds for a smooth hypersurface.  The exact
reference is SGA~2, Expose XII, Corollaire 3.7,
checked in the recomposed primary edition,
\url{https://wstein.org/sga/circle/SGA2.pdf},
printed page 121 (PDF page 129).
An isomorphism must take the positive ample
generator to the positive ample generator.
If the hypersurface degree is $e$, equality of
the hyperplane top degrees would give
\[
                                e=\prod_{i=1}^t a_i.
\]
Adjunction on the two varieties gives, with
the same identified generator,
\[
 e-n-2=\sum_i a_i-(n+t+1),
             \qquad e=1+\sum_i(a_i-1).
\]
Set $u_i=a_i-1\ge0$.  At least two $u_i$
are positive, so expansion of the product
gives
\[
                    \prod_i(1+u_i)>1+\sum_i u_i.
\]
This contradicts the two displayed equalities.
The argument uses intrinsic Picard generators;
it does not assume beforehand that a hypothetical
isomorphism is a projective linear isomorphism.
\end{proof}

\paragraph{Scope, failure boundary and dependencies.}
The ordinary construction may be iterated because
each of its outputs has only isolated singular
special points.  After the deep final step the
special fiber has positive-dimensional intersection
strata; that output is not fed back into the
avoidance lemma.  The result gives complete-height
families for many specified complete-intersection
multidegrees and, in the last theorem, a deep
degeneration with nonzero extreme monodromy.
It does not prove full G for every complete
intersection or allow deformation of an arbitrary
fixed cycle into one of these arithmetic families.

The geometric inputs are the explicitly verified
jet-surjectivity and finite-point avoidance
arguments, elementary field/ideal CRT, projective
complete-intersection cohomology from the Koszul
complex, and the stated primary Picard theorem.
The ordinary local domain uses frozen ON17,
including its arbitrary-thickness and all-new-cycle
extension corollary.  The deep last step uses
CA18 and the separately reviewed RI18.  The
height implication is frozen PH16 and GA17.
Every auxiliary local descent inherits its full
degree and original absolute-value normalization
from these activated interfaces.  No arbitrary
proper actual pushforward, unproved Gysin
cancellation or general resolution-factorization
assumption appears.  General G and arbitrary-
variety complete unconditional height remain open.
